% Generated draft: requires mathematical and source review.
\documentclass[11pt]{article}
\usepackage[margin=25mm]{geometry}
\usepackage{fontspec}
\setmainfont{Latin Modern Roman}
\usepackage{amsmath,amssymb,mathtools,unicode-math}
\setmathfont{Latin Modern Math}
\usepackage{enumitem,needspace,xurl,hyperref}
\setlength{\parindent}{0pt}
\setlength{\parskip}{5pt}
\setlength{\emergencystretch}{3em}
\newcommand{\N}{\mathbb{N}}
\newcommand{\Z}{\mathbb{Z}}
\newcommand{\Q}{\mathbb{Q}}
\newcommand{\R}{\mathbb{R}}
\newcommand{\C}{\mathbb{C}}
\newcommand{\F}{\mathbb{F}}
\newcommand{\A}{\mathbb{A}}
\newcommand{\PP}{\mathbb{P}}
\newcommand{\E}{\mathbb{E}}
\newcommand{\eps}{\varepsilon}
\newcommand{\dd}{\,\mathrm{d}}
\begin{document}
\section*{UnsolvedMath ID: 30004647}
\textbf{Source problem code:} catalogue entry
\subsection{Definitions and mathematical statement}
Let \(X\) be a connected smooth quasi-projective variety over \(\C\), let \(x\in X(\C)\), and let \(X\hookrightarrow\overline X\) be a normal-crossing compactification, with boundary \(D=\overline X\setminus X\). Write \(\pi_1^{\mathrm{top}}(X,x)\) for the fundamental group of the associated complex analytic space. For an integer \(n\geq 1\), a representation
\[
\rho:\pi_1^{\mathrm{top}}(X,x)\longrightarrow\operatorname{GL}_n(\C)
\]
is irreducible if the only \(\rho\)-invariant complex linear subspaces of \(\C^n\) are \(0\) and \(\C^n\). For each irreducible component \(D_i\) of the boundary, let \(T_i\) denote the monodromy of \(\rho\) around a small loop encircling a general smooth point of \(D_i\), transported to the base point. Its conjugacy class is independent of the choices relevant here. Quasi-unipotent local monodromy means that for each \(i\) there exists an integer \(m_i\geq 1\) such that
\[
(T_i^{m_i}-I_n)^n=0.
\]
For a number field \(K\), its ring of integers \(\mathcal O_K\) consists of the elements of \(K\) integral over \(\Z\).

\textbf{Source-attributed conjectural statement (Simpson's integrality conjecture).} For every such \(X,x,\overline X\), every integer \(n\geq 1\), and every rigid irreducible representation \(\rho\) with quasi-unipotent local monodromies, do there exist a number field \(K\), an embedding \(\iota:K\hookrightarrow\C\), and a matrix \(g\in\operatorname{GL}_n(\C)\) such that
\[
g\rho(\gamma)g^{-1}\in\operatorname{GL}_n\bigl(\iota(\mathcal O_K)\bigr)
\qquad\text{for every }\gamma\in\pi_1^{\mathrm{top}}(X,x)?
\]
Here the field, its embedding, and the conjugating matrix may depend on \(\rho\), but must work simultaneously for every \(\gamma\).

\textbf{Unresolved formulation qualification.} The supplied record does not define ``rigid.'' In particular, it does not specify the deformation space, whether local monodromy conjugacy classes or the determinant are held fixed, or whether cohomological rigidity is intended. These choices cannot be silently identified. The displayed question therefore retains the source's unresolved rigidity hypothesis and is not a fully specified conjecture until that hypothesis is clarified. The original statement also refers to local monodromies ``as above,'' but the preceding source passage is not supplied. The boundary-loop description above spells out the usual meaning of the local-monodromy condition in the supplied clean statement; it cannot recover any additional hypotheses from that missing passage.
\subsection{Short English statement}
The supplied statement, attributed to Simpson, asks whether every rigid irreducible complex representation with quasi-unipotent boundary monodromy can be conjugated so that all its matrices lie in the ring of integers of one number field. The meaning of rigidity and any hypotheses hidden in the original phrase ``as above'' require clarification. The supplied open status and review labels are curation claims, not independently verified conclusions.
\subsection{Sources}
\begin{itemize}
\item Local UnsolvedMath registry: \texttt{data/unsolvedmath/problems.json}, permanent ID \(30004647\), title ``Integrality of Rigid Irreducible Local Systems,'' exact source problem code \texttt{catalogue entry}. The original statement calls this ``Simpson's integrality conjecture''; the proposer field is empty, and the supplied proposed year is \(2021\).
\item Supplied source citation: \emph{Homotopic and Geometric Galois Theory} (2022), Oberwolfach Reports, pp.\ 663--744, DOI \texttt{10.4171/owr/2021/12}. Supplied source URL and literature reference labelled ``Oberwolfach Report 12/2021'': \url{https://doi.org/10.4171/owr/2021/12}. The background describes Oberwolfach Reports as records of mathematical workshops and their open-problem discussions. No precise page or passage for this conjecture is supplied, and the cited report has not been consulted here.
\item Supplied literature assessment, checked on \texttt{2026-08-22}: status ``open.'' This assessment says the record is a citation fragment beginning with ``conjecture'' that supplies no complete claim, and that its missing statement cannot reliably be resolved. This is supplied dated curation, not independently verified literature assessment; the background explicitly requests independent verification of time-sensitive status and cited claims. Its description conflicts with the complete-looking clean statement supplied in the same record.
\item The supplied statement status is \texttt{corrected\_verified}, accompanied by a claim of manual semantic, typographic, and LaTeX review establishing self-containment and fidelity. That review claim is preserved as source metadata and has not been independently verified here. The source citation gives the year \(2022\), while the report label, DOI, and proposed-year field refer to \(2021\); the supplied evidence does not explain this dating distinction.
\end{itemize}
\end{document}
